\section{UNION, INTERSECTION, PRODUCT OF A FAMILY OF SETS}

\begin{enumerate}
\def\labelenumi{\arabic{enumi}.}
\tightlist
\item
  In this section we consider a family \((\mathbf{X}_{i})_{i\in\mathbf{I}}\) of subsets of a set E, in which the index set I is arbitrary; the set of subsets of E belonging to the family will be denoted by \(\mathfrak{F}\) (which is therefore a subset of \(\mathfrak{P}(\mathbf{E})\) ).
\end{enumerate}

If I is finite, the consideration of the family \((\mathbf{X}_{t})\) reduces to that of several subsets of E, which may or may not be distinct, and in number equal to the number of elements of I. For example, any three subsets \(X_{1}, X_{2}, X_{3}\) of E form a family of subsets of E, the index set here consisting of the numbers 1, 2, 3.

\begin{enumerate}
\def\labelenumi{\arabic{enumi}.}
\setcounter{enumi}{1}
\tightlist
\item
  Let \(J\) be any subset of \(I\) , and consider the set of all elements \(x\) which have the property ``there exists \(\iota \in J\) such that \(x \in X_\iota\)''. This set is called the union of the family of sets \((X_\iota)_{\iota \in J}\) , and is denoted by \(\bigcup_{\iota\in J} X_\iota\) .
\end{enumerate}

We may also formulate this definition as follows: to the mapping \(\iota \to \mathbf{X}_{\iota}\) of \(\mathbf{I}\) into \(\mathfrak{P}(\mathbf{E})\) there corresponds a well-defined subset \(\mathbf{C}\) of \(\mathbf{I} \times \mathbf{E}\) such that \(\mathbf{X}_{\iota} = \mathbf{C}(\iota)\) ( \(\S 3\) , no. 7), and we have \(\bigcup_{\iota \in J} \mathbf{X}_{\iota} = \mathbf{C}(\mathbf{J})\) .

In particular,

\[
\bigcup_ {t \in \varnothing} \mathrm{X} _ {t} = \mathrm{C} (\varnothing) = \varnothing .\tag{33}
\]

When \(J = I\) , we often write \(\bigcup_{i} X_i\) , or simply \(\bigcup X_i\) , in place of \(\bigcup_{i \in I} X_i\) .

The union \(\bigcup_{i} X_{i}\) depends only on the set \(\mathfrak{F}\) ; in other words, it is the same for two families corresponding to the same subset \(\mathfrak{F}\) of \(\mathfrak{P}(E)\) . In particular, it is equal to the union of the family defined by the canonical mapping of \(\mathfrak{F}\) into \(\mathfrak{P}(E)\) , and we may therefore write \(\bigcup_{X \in \mathfrak{F}} X\) , which is also called the union of the sets belonging to \(\mathfrak{F}\) .

When \(\mathbf{I}\) is a set whose elements are explicitly designated, for example the numbers 1, 2, 3, we have

\[
\bigcup_ {i \in I} X _ {i} = X _ {1} \cup X _ {2} \cup X _ {3},
\]

which justifies the name ``union'' given in the general case to the set \(\bigcup_{i}X_{i}\) .

\begin{enumerate}
\def\labelenumi{\arabic{enumi}.}
\setcounter{enumi}{2}
\tightlist
\item
  For any \(J \subset I\) we have
\end{enumerate}

\[
\bigcup_ {i \in J} X _ {i} \subset \bigcup_ {i \in I} X _ {i}.
\]

In particular, for all \(x \in I\) , we have

\[
\mathrm{X} _ {x} \subset \bigcup_ {i \in I} \mathrm{X} _ {i}.
\]

Conversely, if Y is a subset of E such that \(X_{i} \subset Y\) for all \(i \in I\) , then

\[
\bigcup_ {i} \mathrm{X} _ {i} \subset \mathrm{Y}.
\]

More generally, if \((\mathbf{Y}_{i})\) is another family of subsets of E, indexed by the same set I, and if \(X_{i} \subset Y_{i}\) for all \(i \in I\) , then

\[
\bigcup_ {i} X _ {i} \subset \bigcup_ {i} Y _ {i}.
\]

Let F be another set, and let \(\mathbf{X} \to \mathbf{K}(\mathbf{X})\) be the mapping of \(\mathfrak{P}(\mathbf{E})\) into \(\mathfrak{P}(\mathbf{F})\) defined by a subset K of \(E \times F\) . Then we have

\[
\mathrm{K} \left(\bigcup_ {i \in J} \mathrm{X} _ {i}\right) = \bigcup_ {i \in J} \mathrm{K} (\mathrm{X} _ {i}).\tag{34}
\]

Now let \(\mathbf{L}\) be another index set, and \((J_{\lambda})_{\lambda \in \mathbf{L}}\) a family of subsets of \(\mathbf{I}\) . Then

\[
\bigcup_ {\iota \in \bigcup_ {\lambda \in L} J _ {\lambda}} X _ {\iota} = \bigcup_ {\lambda \in L} \left(\bigcup_ {\iota \in J _ {\lambda}} X _ {\iota}\right).\tag{35}
\]

This is the general associativity formula for unions. When I and L are sets whose elements are explicitly designated, the relations we obtain have already been given (see § 1, no. 14). If L alone satisfies this condition and consists, say, of the numbers 1, 2, we have

\[
\bigcup_ {i \in J _ {1} \cup J _ {2}} X _ {i} = \left(\bigcup_ {i \in J _ {1}} X _ {i}\right) \cup \left(\bigcup_ {i \in J _ {2}} X _ {i}\right).\tag{36}
\]

Let \((\mathbf{X}_i)_{i\in I}\) and \((\mathbf{Y}_x)_{x\in K}\) be any two families of subsets of E. Then we have

\[
\left(\bigcup_ {i \in I} X _ {i}\right) \cap \left(\bigcup_ {x \in K} Y _ {x}\right) = \bigcup_ {(i, x) \in I \times K} (X _ {i} \cap Y _ {x}),\tag{37}
\]

the distributivity formula, which includes the second formula of (10) as a particular case.

If \((\mathbf{X}_t)_{i\in \mathbf{I}}\) is a family of subsets of \(\mathbf{E}\) , and if \((\mathbf{Y}_{\mathbf{x}})_{\mathbf{x}\in \mathbf{K}}\) is a family of subsets of \(\mathbf{F}\) , then

\[
\left(\bigcup_ {i \in I} X _ {i}\right) \times \left(\bigcup_ {x \in K} Y _ {x}\right) = \bigcup_ {(i, x) \in I \times K} (X _ {i} \times Y _ {x}).\tag{38}
\]

\begin{enumerate}
\def\labelenumi{\arabic{enumi}.}
\setcounter{enumi}{3}
\tightlist
\item
  A family \((\mathbf{X}_i)_{i\in I}\) of subsets of \(E\) is a covering of a subset \(A\) of \(E\) , or covers \(A\) , if
\end{enumerate}

\[
\mathrm{A} \subset \bigcup_ {i \in \mathrm{I}} \mathrm{X} _ {i}.
\]

In particular, if \((\mathbf{X}_{t})\) is a covering of E, we have

\[
\bigcup_ {i \in I} X _ {i} = E.
\]

A partition of \(\mathbf{E}\) is a covering \((\mathbf{X}_t)\) of \(\mathbf{E}\) such that

\begin{enumerate}
\def\labelenumi{(\alph{enumi})}
\item
  \(\mathbf{X}_{\iota} \neq \varnothing\) for all \(\iota \in \mathbf{I}\) ;
\item
  \(X_{t} \cap X_{x} = \varnothing\) for each pair of different indices t, x in I. (This second condition may be expressed by saying that the \(X_{t}\) are pairwise disjoint.)
\end{enumerate}

These conditions imply that \(\iota \to X_{\iota}\) is a bijection of I onto the set \(\mathfrak{F}\) of subsets of the partition. Hence, if \(\mathfrak{F}\) is given, the family is determined to within a one-to-one correspondence of index sets. In particular, a partition may be considered indifferently as a set of subsets or as a family of subsets.

\begin{enumerate}
\def\labelenumi{\arabic{enumi}.}
\setcounter{enumi}{4}
\tightlist
\item
  Let \((\mathbf{X}_t)_{i \in I}\) be any family of non-empty subsets of a set \(E\) . In the product \(I \times E\) , let \(\mathbf{X}_t'\) denote the subset \(\{\iota\} \times \mathbf{X}_t\) for each \(\iota \in I\) . The set
\end{enumerate}

\[
S = \bigcup_ {i \in I} X _ {i} ^ {\prime}
\]

is called the sum of the family \((\mathbf{X}_i)_{i\in I}\) . It is clear that the family \((\mathbf{X}_i')_{i\in I}\) is a partition of \(S\) and that for each \(i\in I\) the mapping \(x_i\to (i,x_i)\) is a bijection of \(\mathbf{X}_i\) onto \(\mathbf{X}_i'\) . By abuse of language, any set in one-to-one correspondence with \(S\) is often called the sum of the family \((\mathbf{X}_i)_{i\in I}\) , and the \(\mathbf{X}_i\) are usually identified with the subsets of this set to which they correspond.

The sum of two non-empty sets E and F is often said to be obtained by adjoining the set F to E.

\begin{enumerate}
\def\labelenumi{\arabic{enumi}.}
\setcounter{enumi}{5}
\tightlist
\item
  With the notation of no. 2, the set of elements x of E which have the property ``for all \(\iota\in J\) , \(x\in X_{\iota}\) '' is called the intersection of the family of sets \((\mathbf{X}_i)_{i \in J}\) , and is denoted by \(\bigcap_{i \in J} \mathbf{X}_i\) ; when \(J = I\) , we often write \(\bigcap_{i} \mathbf{X}_i\) , or simply \(\bigcap \mathbf{X}_i\) , instead of \(\bigcap_{i \in I} \mathbf{X}_i\) .
\end{enumerate}

We have

\[
\mathbb {C} \left(\bigcup_ {i \in J} X _ {i}\right) = \bigcap_ {i \in J} \left(\mathbb {C} X _ {i}\right).\tag{39}
\]

In particular, if \(\mathrm{J} = \emptyset\) ,

\[
\bigcap_ {i \in \emptyset} X _ {i} = E.\tag{40}
\]

The intersection \(\bigcap_{i} X_i\) depends only on the set \(\mathfrak{F}\) , and may be written \(\bigcap_{X \in \mathfrak{F}} X\) . For example, if \(I\) consists of the numbers 1, 2, 3, we have

\[
\bigcap_ {i} \mathrm{X} _ {i} = \mathrm{X} _ {1} \cap \mathrm{X} _ {2} \cap \mathrm{X} _ {3}.
\]

\begin{enumerate}
\def\labelenumi{\arabic{enumi}.}
\setcounter{enumi}{6}
\tightlist
\item
  Formula (39) allows us to generalize the duality rule. If a subset A of E is obtained from other subsets X, Y, Z and families \((\mathbf{X}_{t}), (\mathbf{Y}_{x}), (\mathbf{Z}_{\lambda})\) of subsets of E by applying (in any order) only the operations \(\mathbb C, \cup, \cap, \bigcup, \bigcap\), then we shall obtain the complement \(\mathbb C A\) by replacing the subsets X, Y, Z, \(X_{t}, Y_{x}, Z_{\lambda}\) by their complements, and the operations \(\cup, \cap, \bigcup, \bigcap\) by \(\cap, \cup, \bigcap, \bigcup\), respectively, the order of operations being preserved; of course, the operations of intersection and union are not to be altered where they apply to index sets written under the signs \(\bigcup\) and \(\bigcap\).
\end{enumerate}

As in §1, no. 15, we define the dual of a relation \(\mathbf{A} = \mathbf{B}\) or \(\mathbf{A} \subset \mathbf{B}\) where \(\mathbf{A}\) and \(\mathbf{B}\) are subsets of \(\mathbf{E}\) of the above form.

\begin{enumerate}
\def\labelenumi{\arabic{enumi}.}
\setcounter{enumi}{7}
\tightlist
\item
  For all \(J \subset I\) we have
\end{enumerate}

\[
\bigcap_ {i \in I} X _ {i} \subset \bigcap_ {i \in J} X _ {i}.
\]

In particular, for all \(x \in I\) we have

\[
\bigcap_ {i} \mathrm{X} _ {i} \subset \mathrm{X} _ {x}.
\]

Conversely, if \(\mathbf{Y} \subset \mathbf{X}_i\) for all \(\iota \in I\) , then

\[
\mathrm{Y} \subset \bigcap_ {i} \mathrm{X} _ {i}.
\]

More generally, if \((\mathbf{Y}_{t})\) is another family of subsets of E, indexed by the same set I, and if \(X_{t} \subset Y_{t}\) for all \(t \in I\) , then

\[
\bigcap_ {i} X _ {i} \subset \bigcap_ {i} Y _ {i}.
\]

The union of the sets \(X_{t}\) is the intersection of all sets Y such that \(X_{t} \subset Y\) for all \(t \in I\) . The intersection of the \(X_{t}\) is the union of all Z such that \(Z \subset X_{t}\) for all \(t \in I\) .

The following formulae are the duals of (35) and (37), respectively:

\begin{enumerate}
\def\labelenumi{(\arabic{enumi})}
\setcounter{enumi}{40}
\tightlist
\item
\end{enumerate}

\[
\bigcap_ {i \in \bigcup_ {\lambda \in L} J _ {\lambda}} X _ {i} = \bigcap_ {\lambda \in L} \left(\bigcap_ {i \in J _ {\lambda}} X _ {i}\right) \quad (\text { associativity });\tag{42}
\]

\[
\left(\bigcap_ {t \in I} X _ {t}\right) \cup \left(\bigcap_ {x \in K} Y _ {x}\right) = \bigcap_ {(t, x) \in I \times K} (X _ {t} \cup Y _ {x}) \quad (\text { distributivity }).
\]

If \((\mathbf{X}_i)_{i\in I}\) is a family of subsets of \(\mathbf{E}\) , and \((\mathbf{Y}_x)_{x\in \mathbb{K}}\) a family of subset of \(\mathbf{F}\) , then

\[
\left(\bigcap_ {i \in I} X _ {i}\right) \times \left(\bigcap_ {x \in K} Y _ {x}\right) = \bigcap_ {(i, x) \in I \times K} (X _ {i} \times Y _ {x}).\tag{43}
\]

Moreover, if \((\mathbf{X}_{t})\) and \((\mathbf{Y}_{t})\) are families of subsets of E and F, respectively, indexed by the same set I, then

\[
\left(\bigcap_ {i \in I} X _ {i}\right) \times \left(\bigcap_ {i \in I} Y _ {i}\right) = \bigcap_ {i \in I} (X _ {i} \times Y _ {i}).\tag{44}
\]

The formula (34) has no dual; in general all we can say is

\[
\mathrm{K} \left(\bigcap_ {i \in J} X _ {i}\right) \subset \bigcap_ {i \in J} \mathrm{K} (X _ {i}).\tag{45}
\]

We have equality in (45) for all families \((\mathbf{X}_{t})\) only if \(\mathbf{X} \to \mathbf{K}(\mathbf{X})\) is the inverse extension of a mapping of a subset of F into E. Consequently, if f is a mapping of F into E, we have (generalizing (14))

\[
\bar {f} ^ {1} \left(\bigcap_ {i \in J} X _ {i}\right) = \bigcap_ {i \in J} \bar {f} ^ {1} (X _ {i}).\tag{46}
\]

\begin{enumerate}
\def\labelenumi{\arabic{enumi}.}
\setcounter{enumi}{8}
\tightlist
\item
  Let E be any set and let I be any index set. The set of all families \((x_{i})_{i\in I}\) of elements of E, indexed by I, is denoted by \(E^{I}\) , and the operation of passing from E to \(E^{I}\) is called exponentiation. \(E^{I}\) is thus in one-to-one correspondence with the set of all mappings of I into E (which for this reason is often denoted by \(E^{I}\) , by abuse of language), as well as with a subset of \(\mathfrak{P}(\mathbf{I} \times \mathbf{E})\) , by considering the graphs of these mappings. The sets \(E^{J}\) corresponding to subsets J of the set I may therefore be considered all as subsets of the same set, which is in one-to-one correspondence with a subset of \(\mathfrak{P}(\mathbf{I} \times \mathbf{E})\) .
\end{enumerate}

Now let \((\mathbf{X}_i)_{i\in I}\) be a family of subsets of E, indexed by the same set I, and let J be any subset of I. The property ``for all \(i\in J\) , \(x_i\in X_i\)'' of the family \((x_i)_{i\in J}\) defines a subset of \(\mathbf{E}^{\mathbf{J}}\) , called the product of the family of sets \((\mathbf{X}_i)_{i\in J}\) , and denoted by \(\prod_{i\in J}\mathbf{X}_i\) (or simply \(\prod_{i}\mathbf{X}_i\) when \(J = I\) ). The \(\mathbf{X}_i\) are called the factors of the product. Note that \(\prod_{i\in \emptyset}\mathbf{X}_i\) is a set consisting of one element (corresponding to the empty subset of \(\mathbf{I}\times \mathbf{E}\) ). If \(\mathbf{X}_i = \mathbf{E}\) for all \(i\in J\) , then we have

\[
\prod_ {i \in J} X _ {i} = E ^ {J}.
\]

If, for example, I consists of the three numbers 1, 2, 3, then \(\prod_{i\in J}X_i\) is in one-to-one correspondence with the set \(X_1\times X_2\times X_3\) . 10. If \(R\{x,y\}\) is a relation between a generic element \(x\) of a set E and a generic element \(y\) of a set F, then the following propositions are equivalent:

``for each x there exists y such that R\{x,y\}''

and

``there exists a mapping f of E into F such that R\{x, f(x)\} for all x''.

The assertion of this equivalence is known as the axiom of choice (or Zermelo's axiom). We shall sometimes indicate whether the proof of a theorem depends on it.

The axiom of choice is equivalent to the following proposition: ``if, for each \(\iota\in I\) , we have \(X_{\iota}\neq\phi\) , then \(\prod_{\iota\in I}X_{\iota}\neq\phi\) ''.

\begin{enumerate}
\def\labelenumi{\arabic{enumi}.}
\setcounter{enumi}{10}
\tightlist
\item
  In this and the following subsection, we shall consider a non-empty product \(\prod A_{i}\) , where \((A_{i})\) is any family of (non-empty) subsets of E.
\end{enumerate}

Let \(J\) be a subset of \(I\) . The mapping \((x_{i})_{i\in I}\to(x_{i})_{i\in J}\) of \(\prod_{i\in I}A_{i}\) onto \(\prod_{i\in J}A_{i}\) is called the projection of \(\prod_{i\in I}A_{i}\) onto \(\prod_{i\in J}A_{i}\) , and is denoted by \(pr_{J}\) . In particular, the mapping \((x_{i})_{i\in I}\to x_{x}\) of \(\prod_{i\in I}A_{i}\) onto \(A_{x}\) is called the coordinate function (or projection) of index \(x\) , and is denoted by \(pr_{x}\) .

If \(z\) is an element of \(\prod_{i} A_{i}\) , we have \(z = (\mathrm{pr}_{i}z)_{i \in I}\) .

Let \(J_1, J_2\) be two sets forming a partition of I. Then \(z \to (\mathrm{pr}_{J_1}z, \mathrm{pr}_{J_2}z)\) is a one-to-one mapping of \(\prod_{i\in I} A_i\) onto \(\prod_{i\in J_1} A_i \times \prod_{i\in J_2} A_i\) .

In general, if \((\mathbf{J}_{\lambda})_{\lambda \in \mathbf{L}}\) is any partition of the set I, the mapping \(z \to (\mathrm{pr}_{\mathbf{J}_{\lambda}} z)_{\lambda \in \mathbf{L}}\) is a bijection (called canonical) of \(\prod_{i \in I} A_i\) onto the product \(\prod_{\lambda \in L} \left( \prod_{i \in J_{\lambda}} A_i \right)\) . This may also be expressed by saying that the product of a family of sets is associative.

\begin{enumerate}
\def\labelenumi{\arabic{enumi}.}
\setcounter{enumi}{11}
\tightlist
\item
  The following propositions generalize those of §3, no. 3; \((\mathbf{X}_t)\) , \((\mathbf{Y}_t)\) denote families of subsets of E such that \(\mathbf{X}_t \subset \mathbf{A}_t\) and \(\mathbf{Y}_t \subset \mathbf{A}_t\) for all \(t \in I\) ; Z denotes any subset of \(\prod A_t\) .
\end{enumerate}

\begin{enumerate}
\def\labelenumi{(\alph{enumi})}
\item
  If \(\prod_{i} X_{i} \neq \emptyset\) , the relation `` \(\prod_{i} X_{i} \subset \prod_{i} Y_{i}\) '' is equivalent to ``for all \(i \in I\) , \(X_{i} \subset Y_{i}\) ''.
\item
  We have \(\operatorname{pr}_{\chi}^{-1}(\mathbf{X}_{\chi}) = \prod_{i} \mathbf{Y}_{i}\) , where \(\mathbf{Y}_{\chi} = \mathbf{X}_{\chi}\) and \(\mathbf{Y}_{i} = \mathbf{A}_{i}\) whenever \(i \neq \chi\) . Hence
\end{enumerate}

\[
\prod_ {i \in I} X _ {i} = \bigcap_ {i \in I} \operatorname{pr}_{i}^{-1}(X _ {i}).\tag{47}
\]

\begin{enumerate}
\def\labelenumi{(\alph{enumi})}
\setcounter{enumi}{2}
\tightlist
\item
  If \(\prod_{i} X_i \neq \emptyset\) , then
\end{enumerate}

\[
\operatorname{pr} _ {x} \left(\prod_ {i} X _ {i}\right) = X _ {x}.\tag{48}
\]

\begin{enumerate}
\def\labelenumi{(\alph{enumi})}
\setcounter{enumi}{3}
\tightlist
\item
  For all Z we have
\end{enumerate}

\[
\mathbf {Z} \subset \prod_ {\iota} \operatorname{pr} _ {\iota} (\mathbf {Z}).\tag{49}
\]

\begin{enumerate}
\def\labelenumi{(\alph{enumi})}
\setcounter{enumi}{4}
\tightlist
\item
  Let \((\mathbf{J}_1, \mathbf{J}_2)\) be a partition of I into two sets, let \((a_t)_{t \in J_1}\) be a family of elements of E, and let \((\mathbf{X}_t)_{t \in J_2}\) be a family of subsets of E, such that \(a_t \in \mathbf{A}_t\) for all \(t \in J_1\) , and \(\mathbf{X}_t \subset \mathbf{A}_t\) for all \(t \in J_2\) . Then the product \(\prod_{i \in I} \mathbf{Y}_t\) , where \(\mathbf{Y}_t = \{a_t\}\) when \(t \in J_1\) , and \(\mathbf{Y}_t = \mathbf{X}_t\) when \(t \in J_2\) , can be put in one-to-one correspondence with \(\prod_{i \in J_2} \mathbf{X}_t\) by projecting onto this latter set.
\end{enumerate}

\begin{enumerate}
\def\labelenumi{\arabic{enumi}.}
\setcounter{enumi}{12}
\tightlist
\item
  Let \((\mathbf{A}_i)_{i\in I}\) be a family of subsets of a set \(F\) , and let \(f\) be a mapping of a set \(E\) into the product \(\prod_{i\in I} A_i\) . If we put \(f_i(x) = \text{pr}_i(f(x))\) , then \(f_i\) is a mapping of \(E\) into \(A_i\) , and \(f\) is the mapping \(x \to (f_i(x))\) . Conversely, if for each index \(i \in I\) , \(f_i\) is a mapping of \(E\) into \(A_i\) , then
\end{enumerate}

\[
x \rightarrow (f _ {i} (x))
\]

is a mapping of E into \(\prod_{i\in I}A_i\) , which is denoted by \((f_i)\) (by abuse of language, because this notation already denotes the family of mappings \(f_i\) ). Thus we define a bijection (called canonical) of the set \(\left(\prod_{i\in I}A_i\right)^E\) onto the set \(\prod_{i\in I}(A_i^E)\) .

\begin{enumerate}
\def\labelenumi{\arabic{enumi}.}
\setcounter{enumi}{13}
\item
  Let E, F, G be three sets. For each mapping f of \(F \times G\) into E and for each \(y \in G\) , let \(f_{y}\) denote the partial mapping \(x \to f(x, y)\) of F into E. Then \(y \to f_{y}\) is a mapping of G into \(E^{F}\) . Conversely, for each mapping g of G into \(E^{F}\) there exists a unique mapping f of \(F \times G\) into E such that \(f_{y} = g(y)\) for all \(y \in G\) . Thus we define a bijection (called canonical) of the set \(E^{F \times G}\) onto the set \((E^{F})^{G}\) .
\item
  Let \((\mathbf{A}_t)_{t \in \mathbf{I}}\) be a family of non-empty subsets of a set \(\mathbf{E}\) . For each \(t \in I\) let \(f_t\) be a mapping of \(\mathbf{A}_t\) into a set \(\mathbf{F}\) such that, for each pair of indices \((t, x)\) , \(f_t\) and \(f_x\) agree on \(\mathbf{A}_t \cap \mathbf{A}_x\) . Then, if
\end{enumerate}

\[
\mathrm{A} = \bigcup_ {i \in \mathrm{I}} \mathrm{A} _ {i},
\]

there exists a unique mapping \(f\) of \(\mathbf{A}\) into \(\mathbf{F}\) such that the restriction of \(f\) to each \(\mathbf{A}_i\) is equal to \(f_i\) . In particular, if \(\mathbf{A}_i \cap \mathbf{A}_x = \emptyset\) whenever \(i \neq x\) , we see that sets \(\mathbf{F}^{\mathbf{A}}\) and \(\prod_{i \in I} \mathbf{F}^{\mathbf{A}_i}\) are in one-to-one correspondence (called canonical).

